\newpage
%%%%%%%%%%%%%%%%%%%%%%%%%%%%%%%%%%%%%%%%%%%%%%%%%%%%%%%%%%%%%%%%
%%%%%%%%%%%%%%%%%%%%%%%%%%%%%%%%%%%%%%%%%%%%%%%%%%%%%%%%%%%%%%%%
%%%%%%%%%%%%%%%%%%%%%%%%%% Enunciado %%%%%%%%%%%%%%%%%%%%%%%%%%%

\begin{myblock}
\phantomsection\addcontentsline{toc}{section}{Parte B | RNN}
\section*{Parte B | RNN}

    Usando la misma base de la tarea 1, entrena un modelo por cada arquitectura: 
    CNN, RNN/LSTM/GRU, y un Transformer de Hugging Face con fine-tuning. Reporta
    accuracy y F1, incluye matriz de confusión y análisis de errores característicos.  

\end{myblock}

%%%%%%%%%%%%%%%%%%%%%%%%%%%%%%%%%%%%%%%%%%%%%%%%%%%%%%%%%%%%%%%%
%%%%%%%%%%%%%%%%%%%%%%%%%%%%%%%%%%%%%%%%%%%%%%%%%%%%%%%%%%%%%%%%

%%%%%%%%%%%%%%%%%%%%%%%%%%%%%%%%%%%%%%%%%%%%%%%%%%%%%%%%%%%%%%%%
%%%%%%%%%%%%%%%%%%%%%%%%%%%%%%%%%%%%%%%%%%%%%%%%%%%%%%%%%%%%%%%%

Recordemos que, a diferencia de las CNN, que capturan patrones locales posicionalmente invariantes, las RNN
procesan el texto de manera secuencial, manteniendo un estado oculto que teóricamente puede 
codificar información de todo el contexto previo. Para tareas de clasificación, este enfoque
permite modelar dependencias de largo alcance y el orden de las palabras. 

Una mejora realizada para las RNN base son las RNN Bidireccionales, que procesa la secuencia en ambas 
direcciones, tanto de izquierda a derecha, y de derecha a izquierda. Esto permite que el 
estado oculto en cualquier punto contenga información tanto del contexto pasado como del 
futuro, resultando en representaciones más ricas. 

\subsection{Arquitectura RNN}

Para obtener un único vector que representa toda la secuencia, se aplica una operación de pooling
sobre la secuencia de estados ocultos de la Bi-RNN. Si $H = [h_1, h_2,..., h_L]$ es la
secuencia de estados ocultos concatenados, donde cada $h_t$ combina el estado hacia adelante y 
hacia atrás, la representación final $h_{\text{pooled}}$ se obtiene como:

\[
    h_{\text{pooled}} = \text{Pool}(H)
\]

El vector resultante $h_{\text{pooled}}$ se pasa a una capa lineal final con una función de 
activación Softmax para la clasificación. La función de pérdida en este caso es la Entropía 
Cruzada, podnerada para manejar el desbalance de clases. Recordemos que en generación de texto, 
la métrica más usada de validación es la Perplejidad (PPL), que no es otra cosa más que la exponencial
de la Entropía Cruzada. 

\[
    \hat{y} = \text{Softmax}(W_o \cdot h_{\text{pooled}} + b_o) \;\;\;\;\;\;\;  L = \text{CrossEntropy}(y, \hat{y})
\]

Los modelos con los que experimentamos partieron de una base con capa de Bi-RNN y MaxPooling. al
oservar un sobreajuste significativo, con diferencia porcentuial de hasta 30 puntos de F1 entre 
training y validación, decidimos explorar diversas estrategias de regularización: aumento en 
el dropout, weight decay, label smoothing y diferentes estrategias de pooling como mean y maxpooling+mean, 
para intentar encontrar una mejor generalización.

\subsection{Resultados}

El primer modelo superó a la CNN, pero mostró un claro sobreajuste (Train F1 $\approx 0.78$ vs. Val F1 $\approx 0.47$). Los siguientes experimentos buscaron mejorar la regularización y explorar el impacto del mecanismo de \emph{pooling}.

\begin{table}[h!]
\centering
\caption{Bi-RNN — Resumen de experimentos en MeIA (polaridad 1–5).}
\label{tab:rnn_meia_transposed}
\begin{tabular}{lccc}
\toprule
\textbf{Parámetro} & \textbf{\texttt{...\_brnn\_five (baseline)}} & \texttt{...\_mean\_do60} & \texttt{...\_maxmean\_ls002} \\
\midrule
\textbf{Pooling}           & \textbf{max}          & mean                & max+mean \\
\textbf{Proj. Dropout}     & \textbf{0.5}          & 0.6                 & 0.5 \\
\textbf{Weight Decay}      & \textbf{$1\mathrm{e}{-4}$} & $2\mathrm{e}{-4}$   & $2\mathrm{e}{-4}$ \\
\textbf{Label Smoothing}   & \textbf{0.0}          & 0.05                & 0.02 \\
\textbf{Val F1-M}          & \textbf{0.467}        & 0.430               & 0.451 \\
\textbf{Test Acc / F1-M}   & \textbf{0.492 / 0.484} & 0.462 / 0.444       & 0.464 / 0.440 \\
\bottomrule
\end{tabular}
\end{table}

\subsection{Visualización de Modelos}

% --- Figura 1: Matrices de Confusión ---
\begin{figure}[h!]
    \centering
    % Subfigura para el primer experimento (el mejor)
    \begin{subfigure}[b]{0.32\textwidth}
        \centering
        % Reemplaza con la ruta de tu imagen
        \includegraphics[width=\textwidth]{/home/cesar/project_b_t02nlp/models/rnn/exp_brnn_five/confusion_test.png} 
        \caption{\textbf{\texttt{exp\_brnn\_five}}}
        \label{fig:conf_rnn_exp1}
    \end{subfigure}
    \hfill 
    % Subfigura para el segundo experimento
    \begin{subfigure}[b]{0.32\textwidth}
        \centering
        % Reemplaza con la ruta de tu imagen
        \includegraphics[width=\textwidth]{/home/cesar/project_b_t02nlp/models/rnn/exp_brnn_ls005_wd2e-4_mean_do60/confusion_test.png}
        \caption{\texttt{...\_mean\_do60}}
        \label{fig:conf_rnn_exp2}
    \end{subfigure}
    \hfill
    % Subfigura para el tercer experimento
    \begin{subfigure}[b]{0.32\textwidth}
        \centering
        % Reemplaza con la ruta de tu imagen
        \includegraphics[width=\textwidth]{/home/cesar/project_b_t02nlp/models/rnn/exp_brnn_maxmean_ls002_wd2e-4_do50/confusion_test.png} 
        \caption{\texttt{...\_maxmean\_ls002}}
        \label{fig:conf_rnn_exp3}
    \end{subfigure}
    \caption{Matrices de confusión para los tres experimentos de Bi-RNN en el dataset MeIA. El modelo baseline (izquierda) muestra el mejor balance.}
    \label{fig:confusion_matrices_rnn}
\end{figure}

% --- Figura 2: Curvas ROC OVR ---
\begin{figure}[h!]
    \centering
    % Subfigura para el primer experimento (el mejor)
    \begin{subfigure}[b]{0.32\textwidth}
        \centering
        % Reemplaza con la ruta de tu imagen
        \includegraphics[width=\textwidth]{/home/cesar/project_b_t02nlp/models/rnn/exp_brnn_five/roc_test.png}
        \caption{\textbf{\texttt{exp\_brnn\_five}}}
        \label{fig:roc_rnn_exp1}
    \end{subfigure}
    \hfill
    % Subfigura para el segundo experimento
    \begin{subfigure}[b]{0.32\textwidth}
        \centering
        % Reemplaza con la ruta de tu imagen
        \includegraphics[width=\textwidth]{/home/cesar/project_b_t02nlp/models/rnn/exp_brnn_ls005_wd2e-4_mean_do60/roc_test.png}
        \caption{\texttt{...\_mean\_do60}}
        \label{fig:roc_rnn_exp2}
    \end{subfigure}
    \hfill
    % Subfigura para el tercer experimento
    \begin{subfigure}[b]{0.32\textwidth}
        \centering
        % Reemplaza con la ruta de tu imagen
        \includegraphics[width=\textwidth]{/home/cesar/project_b_t02nlp/models/rnn/exp_brnn_maxmean_ls002_wd2e-4_do50/roc_test.png}
        \caption{\texttt{...\_maxmean\_ls002}}
        \label{fig:roc_rnn_exp3}
    \end{subfigure}
    \caption{Curvas ROC One-vs-Rest (OVR) para los experimentos de Bi-RNN. Evalúan la capacidad de discriminación de cada clase.}
    \label{fig:roc_curves_rnn}
\end{figure}


\subsection{Conclusiones}

\begin{itemize}
    \item La arquitectura Bi-RNN con MaxPooling superó al mejor modelo de mi TextCNN, estableciendo
    una nueva línea base con un F1-macro de 0.484.
    \item El mecanismo de pooling resultó ser el factor más determinante, MaxPooling fue superior
    a MeanPooling y a la concatenación \texttt{max+mean}. Esto nos está indicando que para este problema
    las señales con alta carga semántica son más informativas que la evidencia contextual distribuida por
    el texto. 
    \item Los intentos de regularización más agresiva, como label smoothing, o un mayor weight decay, 
    fueron contraproducentes. Suavizar las predicciones parece diluir las distinciones sutiles entre las
    clases intermedias (2-3-4), empeorando el rendimiento general. 
    \item Al igual que con la CNN, el desafío principal es la zona media de la polaridad, las clases
    intermedias ya mencionadas. Aunque el modelo captura bien los extremos (1 y 5), la naturaleza ordinal
    y la ambiguedad semántica de las puntuaciones intermedias son la principal causa de error.
\end{itemize}

Este modleo basado en una RNN bidireccional demostró ser más efectivo que la convolucional para capturar
el contexto secuencial de las reseñas. La mejor configuración, una Bi-RNN con MaxPooling sobre los 
estados ocultos y una regularización moderada obtuvieron los mejores resultados de las pruebas, con un
accuracy de 0.492, y un F1-Macro de 0.484 en el conjunto de prueba. El análisis de los experimentos
revelaron que la estrategia de Pooling fue la que más eficientemente abordó el problema, el uso de 
MaxPooling logra identificar características secuenciales más destacadas, lo cual fue benéfico para
promediar la evidencia a lo largo de toda la secuencia, que tendía a centralizar las prediciones. 

Pese a las mejoras, este modelo aún comparte el problema de la zona media, con las etiquetas de polaridad
(2-3-4). La Bi-RNN se establece como una línea base sólida, confirmando que el modelado del 
orden secuencial puede aportar información de valor, y siendo un excelente punto de partidad para evaluar
el portencial de arquitecturas basadas en atención como los Transformers. 








\clearpage